\documentclass[12pt,a4paper]{article}
\usepackage[utf8]{inputenc}
\usepackage[T1]{fontenc}
\usepackage[brazilian]{babel}
\usepackage{geometry}
\usepackage{booktabs}
\usepackage{array}
\usepackage{multirow}
\usepackage{enumitem}
\usepackage{amsmath}
\usepackage{xcolor}
\usepackage{colortbl}
\usepackage{titlesec}

\geometry{margin=2.5cm}

\titleformat{\section}{\large\bfseries}{}{0em}{}
\titleformat{\subsection}{\normalsize\bfseries}{}{0em}{}

\title{\textbf{Teste de Caixa Preta --- Problema dos Triângulos}\\[0.3em]
\large Particionamento de Equivalência e Análise do Valor Limite}
\author{Confiabilidade e Segurança de Software}
\date{}

\begin{document}
\maketitle

\section{Especificação}

Uma função recebe três inteiros (\texttt{a}, \texttt{b}, \texttt{c}) representando os lados de um triângulo e retorna:

\begin{itemize}[nosep]
  \item \textbf{Equilátero} --- se os três lados são iguais
  \item \textbf{Isósceles} --- se exatamente dois lados são iguais
  \item \textbf{Escaleno} --- se todos os lados são diferentes
  \item \textbf{Erro} --- se não formam um triângulo válido
\end{itemize}

\noindent\textbf{Regras de validade:}
\begin{itemize}[nosep]
  \item Todos os lados devem ser positivos: $a > 0$, $b > 0$, $c > 0$
  \item Desigualdade triangular: $a + b > c$,\; $a + c > b$,\; $b + c > a$
\end{itemize}

%% ============================================================
\section{Particionamento de Equivalência}
%% ============================================================

\subsection{Classes de Equivalência --- Validade}

\begin{center}
\begin{tabular}{clc}
\toprule
\textbf{Classe} & \textbf{Condição} & \textbf{Tipo} \\
\midrule
I1 & Algum lado $\leq 0$ & Inválida \\
I2a & $a + b \leq c$ & Inválida \\
I2b & $a + c \leq b$ & Inválida \\
I2c & $b + c \leq a$ & Inválida \\
\bottomrule
\end{tabular}
\end{center}

\subsection{Classes de Equivalência --- Classificação (válidas)}

\begin{center}
\begin{tabular}{clc}
\toprule
\textbf{Classe} & \textbf{Condição} & \textbf{Saída} \\
\midrule
V1  & $a = b = c$ & Equilátero \\
V2a & $a = b \neq c$ & Isósceles \\
V2b & $a = c \neq b$ & Isósceles \\
V2c & $b = c \neq a$ & Isósceles \\
V3  & $a \neq b \neq c \neq a$ & Escaleno \\
\bottomrule
\end{tabular}
\end{center}

\subsection{Casos de Teste --- Particionamento}

\begin{center}
\begin{tabular}{cccccc}
\toprule
\textbf{Caso} & \textbf{a} & \textbf{b} & \textbf{c} & \textbf{Classe} & \textbf{Resultado esperado} \\
\midrule
1  & 5  & 5  & 5  & V1  & Equilátero \\
2  & 3  & 3  & 5  & V2a & Isósceles \\
3  & 3  & 5  & 3  & V2b & Isósceles \\
4  & 5  & 3  & 3  & V2c & Isósceles \\
5  & 3  & 4  & 5  & V3  & Escaleno \\
\midrule
6  & 0  & 5  & 5  & I1  & Erro \\
7  & $-1$ & 5  & 5  & I1  & Erro \\
8  & 1  & 2  & 10 & I2a & Erro \\
9  & 1  & 10 & 2  & I2b & Erro \\
10 & 10 & 1  & 2  & I2c & Erro \\
\bottomrule
\end{tabular}
\end{center}

\noindent As classes inválidas são testadas \textbf{isoladamente} (uma entrada inválida por vez) para garantir que cada validação funciona independentemente.

%% ============================================================
\section{Análise do Valor Limite}
%% ============================================================

No problema do triângulo, a fronteira mais importante é a \textbf{desigualdade triangular}: o ponto onde $a + b = c$ (triângulo degenerado, \textbf{não} forma triângulo) versus $a + b > c$ (forma triângulo).

\subsection{Fronteira: $a + b$ vs $c$ \quad (fixando $a = 3$, $b = 4$)}

\begin{center}
\begin{tabular}{cccclc}
\toprule
\textbf{Caso} & \textbf{a} & \textbf{b} & \textbf{c} & \textbf{Condição} & \textbf{Resultado} \\
\midrule
11 & 3 & 4 & 6 & $a + b > c$ (interior) & Escaleno \\
12 & 3 & 4 & 7 & $a + b = c$ (fronteira) & Erro \\
13 & 3 & 4 & 8 & $a + b < c$ (exterior) & Erro \\
\bottomrule
\end{tabular}
\end{center}

\subsection{Fronteira: $a + c$ vs $b$ \quad (fixando $a = 3$, $c = 4$)}

\begin{center}
\begin{tabular}{cccclc}
\toprule
\textbf{Caso} & \textbf{a} & \textbf{b} & \textbf{c} & \textbf{Condição} & \textbf{Resultado} \\
\midrule
14 & 3 & 6 & 4 & $a + c > b$ & Escaleno \\
15 & 3 & 7 & 4 & $a + c = b$ & Erro \\
16 & 3 & 8 & 4 & $a + c < b$ & Erro \\
\bottomrule
\end{tabular}
\end{center}

\subsection{Fronteira: $b + c$ vs $a$ \quad (fixando $b = 3$, $c = 4$)}

\begin{center}
\begin{tabular}{cccclc}
\toprule
\textbf{Caso} & \textbf{a} & \textbf{b} & \textbf{c} & \textbf{Condição} & \textbf{Resultado} \\
\midrule
17 & 6 & 3 & 4 & $b + c > a$ & Escaleno \\
18 & 7 & 3 & 4 & $b + c = a$ & Erro \\
19 & 8 & 3 & 4 & $b + c < a$ & Erro \\
\bottomrule
\end{tabular}
\end{center}

\subsection{Fronteira: lado $= 0$ (validade do lado)}

\begin{center}
\begin{tabular}{ccccl}
\toprule
\textbf{Caso} & \textbf{a} & \textbf{b} & \textbf{c} & \textbf{Resultado} \\
\midrule
20 & 0 & 5 & 5 & Erro \\
21 & 1 & 5 & 5 & Isósceles \\
22 & 5 & 0 & 5 & Erro \\
23 & 5 & 1 & 5 & Isósceles \\
24 & 5 & 5 & 0 & Erro \\
25 & 5 & 5 & 1 & Isósceles \\
\bottomrule
\end{tabular}
\end{center}

\subsection{Fronteira: Isósceles $\leftrightarrow$ Equilátero}

Caso sutil --- quando o triângulo isósceles ``vira'' equilátero:

\begin{center}
\begin{tabular}{ccccl}
\toprule
\textbf{Caso} & \textbf{a} & \textbf{b} & \textbf{c} & \textbf{Resultado} \\
\midrule
26 & 5 & 5 & 4 & Isósceles \\
27 & 5 & 5 & 5 & Equilátero \\
28 & 5 & 5 & 6 & Isósceles \\
\bottomrule
\end{tabular}
\end{center}

%% ============================================================
\section{Resumo}
%% ============================================================

\begin{center}
\begin{tabular}{lll}
\toprule
\textbf{Técnica} & \textbf{Foco} & \textbf{Casos} \\
\midrule
Particionamento & Um representante de cada classe & 10 \\
Valor Limite & Fronteiras da desigualdade, lado $= 0$, transição isósc./equil. & 18 \\
\midrule
\textbf{Total} & & \textbf{28} \\
\bottomrule
\end{tabular}
\end{center}

\noindent\textbf{Bugs típicos detectados por esses testes:}
\begin{itemize}[nosep]
  \item \texttt{>=} em vez de \texttt{>} na desigualdade triangular (casos 12, 15, 18)
  \item Esquecimento de validar lado igual a zero (casos 20, 22, 24)
  \item Não verificar todas as 3 permutações da desigualdade (casos em I2a/b/c)
  \item Classificar equilátero como isósceles ou vice-versa (casos 26--28)
\end{itemize}

\end{document}
